\documentclass[11pt]{article}
\usepackage[margin=1in]{geometry}
\usepackage{amsmath,amssymb,bm}
\usepackage{array,booktabs}
\usepackage[hidelinks]{hyperref}
\title{Part 4: Detailed Base PINODE Formulation}
\author{}
\date{}
\begin{document}
\maketitle

\section*{Part 4 --- Detailed Base PINODE Formulation}

We now define the \textbf{base physics-informed Neural ODE (PINODE)}. This method sits exactly between the data-only NODE of Part 3 and the hard-projected EBP-PINODE of Part 5.

The defining distinction is
\begin{align}
\boxed{\text{NODE learns an unrestricted neural derivative}}
\end{align}
while
\begin{align}
\boxed{\text{base PINODE learns a neural derivative but penalizes violations of the full RC equations softly.}}
\end{align}

No projection is applied. Therefore the RC balance is encouraged, not guaranteed.

\section*{1. Relationship to Parts 1 and 3}

Everything below preserves the previously fixed contracts:
\begin{align}
\boxed{\text{All-to-One,\ Identity/IND,\ Identity/DEP1,\ Identity/DEP2}},
\end{align}
\begin{align}
\boxed{1C\qquad\text{and}\qquad 2C},
\end{align}
the raw forcing vectors \(v\), training-only normalization \(\mu_x,S_x,\mu_v,S_v,S_y\), dimensionless time
\begin{align}
\boxed{\tau=\frac{t-t_{\rm ref}}{\Delta t}},
\end{align}
training rollout length
\begin{align}
\boxed{N_r},
\end{align}
2C encoder-history length
\begin{align}
\boxed{L_e},
\end{align}
and RK4 substeps per 5-minute sampled interval
\begin{align}
\boxed{N_s}.
\end{align}

The new ingredient is
\begin{align}
\boxed{\mathcal L_{\rm phys}}.
\end{align}

\section*{2. General base-PINODE model}

Define normalized state and forcing as
\begin{align}
z&=S_x^{-1}(x-\mu_x),\\
\widetilde v&=S_v^{-1}(v-\mu_v).
\end{align}

The neural vector field is
\begin{align}
\boxed{\frac{dz}{d\tau}=g_\omega(z,\widetilde v).}
\tag{P4.1}
\end{align}

Converting to physical-time derivative,
\begin{align}
\boxed{
\dot x
=
\frac{dx}{dt}
=
\frac{S_x}{\Delta t}
g_\omega(z,\widetilde v).
}
\tag{P4.2}
\end{align}

Define the raw physical neural derivative
\begin{align}
\boxed{
\widetilde f_\omega(x,v)
=
\frac{S_x}{\Delta t}
g_\omega
\left(
S_x^{-1}(x-\mu_x),
S_v^{-1}(v-\mu_v)
\right).
}
\tag{P4.3}
\end{align}

Thus
\begin{align}
\boxed{\dot x=\widetilde f_\omega(x,v).}
\tag{P4.4}
\end{align}

The tilde emphasizes that this is the \textbf{unconstrained neural derivative}. In base PINODE it is integrated without hard projection.

\section*{3. Neural-network architecture and architecture hyperparameters}

Define
\begin{align}
\xi=
\begin{bmatrix}
z\\
\widetilde v
\end{bmatrix}.
\end{align}

Let the number of hidden layers be
\begin{align}
\boxed{L_{\rm hid}\in\mathbb N,}
\tag{P4.5}
\end{align}
and let hidden layer \(\ell\) contain
\begin{align}
\boxed{n_{\rm hid}^{(\ell)}}
\tag{P4.6}
\end{align}
neurons.

The hidden-width vector is
\begin{align}
\boxed{
\bm n_{\rm hid}
=
\left[
n_{\rm hid}^{(1)},\ldots,n_{\rm hid}^{(L_{\rm hid})}
\right].
}
\tag{P4.7}
\end{align}

Set
\begin{align}
\boxed{h^{(0)}=\xi.}
\tag{P4.8}
\end{align}

For
\begin{align}
\ell=1,\ldots,L_{\rm hid},
\end{align}
define
\begin{align}
\boxed{
h^{(\ell)}
=
\varphi
\left(
W^{(\ell)}h^{(\ell-1)}+b^{(\ell)}
\right).
}
\tag{P4.9}
\end{align}

The output layer is
\begin{align}
\boxed{
g_\omega(z,\widetilde v)
=
W^{(L_{\rm hid}+1)}h^{(L_{\rm hid})}
+
b^{(L_{\rm hid}+1)}.
}
\tag{P4.10}
\end{align}

Its output dimension equals the state dimension:
\begin{align}
\boxed{\dim g_\omega=n_x.}
\tag{P4.11}
\end{align}

The following are explicitly hyperparameters:
\begin{align}
\boxed{L_{\rm hid}=\text{number of hidden layers},}
\tag{P4.12}
\end{align}
and
\begin{align}
\boxed{\bm n_{\rm hid}=\text{number of neurons in the hidden layers}.}
\tag{P4.13}
\end{align}

They are selected during the smaller representative-dataset tuning stage and frozen before full-data training. They are not retuned using final test performance.

\section*{4. Physical RC parameters are jointly trainable}

Base PINODE also contains a physical parameter vector
\begin{align}
\boxed{\theta_{RC}.}
\end{align}

For All-to-One 1C,
\begin{align}
\boxed{
\theta_A^{1C}
=
\{R_{Ao},C_A,\eta_{r,A}\}.
}
\tag{P4.14}
\end{align}

For DEP1 2C,
\begin{align}
\boxed{
\begin{aligned}
\theta_{\rm DEP1}^{2C}
=
\{&
R_{Do},R_{Ko},R_{DK},
R_{Dm},R_{Km},\\
&
C_{a,D},C_{m,D},
C_{a,K},C_{m,K},
\eta_{r,D},\eta_{r,K}
\}.
\end{aligned}
}
\tag{P4.15}
\end{align}

For DEP2 the model additionally uses \(\alpha_c,\alpha_r\), which generate constrained disturbance-allocation factors.

The fairness rule is
\begin{align}
\boxed{\text{base PINODE learns its own RC parameters jointly with the neural vector field.}}
\tag{P4.16}
\end{align}

The inverse-PINN estimates are not handed to base PINODE as truth.

\section*{5. Physical parameter transformations}

Every resistance and capacitance remains positive:
\begin{align}
\boxed{
R_j=R_{j,\min}+\operatorname{softplus}(\rho_{R_j}),
}
\tag{P4.17}
\end{align}
\begin{align}
\boxed{
C_j=C_{j,\min}+\operatorname{softplus}(\rho_{C_j}).
}
\tag{P4.18}
\end{align}

For learnable radiative participation,
\begin{align}
\boxed{\eta_{r,j}=\sigma(\rho_{\eta_j}).}
\tag{P4.19}
\end{align}

For DEP2,
\begin{align}
\boxed{\lambda_{c,D}=2\sigma(\alpha_c),}
\tag{P4.20}
\end{align}
\begin{align}
\boxed{\lambda_{c,K}=2[1-\sigma(\alpha_c)],}
\tag{P4.21}
\end{align}
\begin{align}
\boxed{\lambda_{r,D}=2\sigma(\alpha_r),}
\tag{P4.22}
\end{align}
and
\begin{align}
\boxed{\lambda_{r,K}=2[1-\sigma(\alpha_r)].}
\tag{P4.23}
\end{align}

Therefore
\begin{align}
\boxed{\lambda_{c,D}+\lambda_{c,K}=2,}
\tag{P4.24}
\end{align}
and
\begin{align}
\boxed{\lambda_{r,D}+\lambda_{r,K}=2.}
\tag{P4.25}
\end{align}

\section*{6. Meaning of \(\alpha_r\), \(\lambda_r\), and \(\eta_r\)}

The variable
\begin{align}
\boxed{\alpha_r\in\mathbb R}
\tag{P4.26}
\end{align}
is an unconstrained trainable scalar used to generate the two constrained DEP2 radiative allocation factors.

If
\begin{align}
\alpha_r=0,
\end{align}
then
\begin{align}
\sigma(0)=\frac12,
\end{align}
and therefore
\begin{align}
\boxed{\lambda_{r,D}=\lambda_{r,K}=1.}
\tag{P4.27}
\end{align}

If, for example,
\begin{align}
\sigma(\alpha_r)=0.7,
\end{align}
then
\begin{align}
\boxed{
\lambda_{r,D}=1.4,
\qquad
\lambda_{r,K}=0.6.
}
\tag{P4.28}
\end{align}

Thus \(\alpha_r\) is an optimization variable, while \(\lambda_r\) is the corresponding constrained between-zone radiative allocation.

For DEP2 2C Dining,
\begin{align}
Q_{r,D}^{DEP2}
=
\lambda_{r,D}\bar Q_r.
\end{align}

Mass receives
\begin{align}
\boxed{
Q_{r,D\rightarrow m}
=
\eta_{r,D}\lambda_{r,D}\bar Q_r,
}
\tag{P4.29}
\end{align}
and air receives
\begin{align}
\boxed{
Q_{r,D\rightarrow a}
=
(1-\eta_{r,D})\lambda_{r,D}\bar Q_r.
}
\tag{P4.30}
\end{align}

Therefore
\begin{align}
\boxed{\lambda_r=\text{between-zone allocation},}
\tag{P4.31}
\end{align}
whereas
\begin{align}
\boxed{\eta_r=\text{within-zone air/mass allocation}.}
\tag{P4.32}
\end{align}

\section*{7. When \(\lambda_r\) and \(\eta_r\) cannot be learned together}

For DEP2 1C the radiative term contains
\begin{align}
\eta_r\lambda_r\bar Q_r.
\end{align}

Suppose the required effective multiplier is \(0.6\). Then
\begin{align}
0.6\times1.0&=0.6,\\
0.75\times0.8&=0.6,\\
0.5\times1.2&=0.6.
\end{align}

Thus the 1C equation primarily sees
\begin{align}
\boxed{\eta_r\lambda_r}
\tag{P4.33}
\end{align}
rather than the two parameters separately.

Therefore the primary DEP2 1C model uses
\begin{align}
\boxed{\eta_{r,D}=\eta_{r,K}=1,}
\tag{P4.34}
\end{align}
and learns \(\alpha_r\), hence \(\lambda_r\).

We lock
\begin{align}
\boxed{
\text{do not jointly learn }\eta_r\text{ and }\lambda_r
\text{ in the primary DEP2 1C fit.}
}
\tag{P4.35}
\end{align}

\section*{8. When \(\lambda_r\) and \(\eta_r\) can be learned together}

For DEP2 2C, air receives
\begin{align}
(1-\eta_r)\lambda_r\bar Q_r,
\end{align}
while mass receives
\begin{align}
\eta_r\lambda_r\bar Q_r.
\end{align}

Their sum is
\begin{align}
\boxed{
(1-\eta_r)\lambda_r\bar Q_r
+
\eta_r\lambda_r\bar Q_r
=
\lambda_r\bar Q_r.
}
\tag{P4.36}
\end{align}

For example, if
\begin{align}
\lambda_r=1.2,
\qquad
\eta_r=0.75,
\end{align}
then
\begin{align}
Q_{mass}&=0.9\bar Q_r,\\
Q_{air}&=0.3\bar Q_r.
\end{align}

Hence
\begin{align}
\boxed{Q_{air}+Q_{mass}=1.2\bar Q_r,}
\tag{P4.37}
\end{align}
while
\begin{align}
\boxed{
\frac{Q_{mass}}{Q_{air}+Q_{mass}}
=
0.75
=
\eta_r.
}
\tag{P4.38}
\end{align}

Therefore the exact 1C multiplicative degeneracy is removed. This does not prove perfect statistical identifiability, but it supports jointly estimating
\begin{align}
\boxed{
\eta_{r,D},\eta_{r,K},\alpha_r
}
\tag{P4.39}
\end{align}
for primary DEP2 2C.

\section*{9. What ``physics-informed'' means}

Write the Part-1 RC equations abstractly as
\begin{align}
\boxed{
M_\theta\dot x=b_\theta(x,v).
}
\tag{P4.40}
\end{align}

For the neural derivative,
\begin{align}
\dot x=\widetilde f_\omega(x,v),
\end{align}
define
\begin{align}
\boxed{
r_{\rm RC}(x,v)
=
M_\theta\widetilde f_\omega(x,v)
-
b_\theta(x,v).
}
\tag{P4.41}
\end{align}

The physics loss penalizes
\begin{align}
\boxed{\mathcal L_{\rm phys}\propto\|r_{\rm RC}\|^2.}
\tag{P4.42}
\end{align}

But
\begin{align}
\boxed{r_{\rm RC}\neq0\text{ is allowed during and after training.}}
\tag{P4.43}
\end{align}

Thus base PINODE is a \textbf{soft-physics} method.

\section*{10. Why the full RC residual is used}

For 2C one zone,
\begin{align}
\boxed{
r_{\rm RC}
=
\begin{bmatrix}
r_a\\
r_m
\end{bmatrix}.
}
\tag{P4.44}
\end{align}

For coupled DEP1/DEP2 2C,
\begin{align}
\boxed{
r_{\rm RC}
=
\begin{bmatrix}
r_{a,D}\\
r_{m,D}\\
r_{a,K}\\
r_{m,K}
\end{bmatrix}.
}
\tag{P4.45}
\end{align}

The base PINODE therefore softly penalizes the complete chosen RC model, not merely the total-energy balance.

\section*{11. Where the physics residual is evaluated}

The physics residual is evaluated on model-generated states.

At an internal solver state \(x^{(q)}\) with held forcing \(v_k\),
\begin{align}
\boxed{
r_{\rm RC}^{(q)}
=
M_\theta
\widetilde f_\omega(x^{(q)},v_k)
-
b_\theta(x^{(q)},v_k).
}
\tag{P4.46}
\end{align}

\section*{12. Physics-loss indexing and the four summations}

The generic physics loss is
\begin{align}
\boxed{
\mathcal L_{\rm phys}
=
\frac{1}
{
|\mathcal W_{\rm tr}|N_rN_s4
}
\sum_w
\sum_{\ell=0}^{N_r-1}
\sum_{n=1}^{N_s}
\sum_{q=1}^{4}
\left\|
D^{-1}
r_{\rm RC}^{w,\ell,n,q}
\right\|_2^2.
}
\tag{P4.47}
\end{align}

The four indices are
\begin{align}
\boxed{w=\text{training-rollout index},}
\tag{P4.48}
\end{align}
\begin{align}
\boxed{\ell=\text{5-minute sampled interval within the rollout},}
\tag{P4.49}
\end{align}
\begin{align}
\boxed{n=\text{RK4 numerical substep within that interval},}
\tag{P4.50}
\end{align}
and
\begin{align}
\boxed{q=\text{RK4 stage within that substep}.}
\tag{P4.51}
\end{align}

If \(N_r=6\), then
\begin{align}
\ell=0,\ldots,5
\end{align}
indexes six recursive 5-minute transitions.

If \(N_s=5\), then each 5-minute interval is divided into five solver substeps with
\begin{align}
\boxed{
\Delta t_{\rm sub}
=
\frac{300}{N_s}\ {\rm s}
=
60\ {\rm s}.
}
\tag{P4.52}
\end{align}

Each RK4 substep evaluates the vector field four times, giving \(q=1,\ldots,4\).

For example, if
\begin{align}
|\mathcal W_{\rm tr}|=100,\qquad N_r=6,\qquad N_s=5,
\end{align}
then the residual is evaluated
\begin{align}
\boxed{
100\times6\times5\times4=12{,}000
}
\tag{P4.53}
\end{align}
times.

Hence
\begin{align}
\boxed{
\mathcal L_{\rm phys}
=
\text{mean squared RC violation over all neural derivative evaluations during training}.
}
\tag{P4.54}
\end{align}

\section*{13. Why evaluate physics at RK4 stage points}

The ODE solver visits intermediate states between measured samples. Physics is therefore checked wherever the neural vector field is evaluated.

For normalized state \(z_n\) and held forcing \(\widetilde v_k\),
\begin{align}
k_1&=g_\omega(z_n,\widetilde v_k),\\
k_2&=g_\omega\left(z_n+\frac{h}{2}k_1,\widetilde v_k\right),\\
k_3&=g_\omega\left(z_n+\frac{h}{2}k_2,\widetilde v_k\right),\\
k_4&=g_\omega\left(z_n+hk_3,\widetilde v_k\right).
\end{align}

The normalized stage states are
\begin{align}
\boxed{z_n^{(1)}=z_n,}
\tag{P4.55}
\end{align}
\begin{align}
\boxed{z_n^{(2)}=z_n+\frac{h}{2}k_1,}
\tag{P4.56}
\end{align}
\begin{align}
\boxed{z_n^{(3)}=z_n+\frac{h}{2}k_2,}
\tag{P4.57}
\end{align}
and
\begin{align}
\boxed{z_n^{(4)}=z_n+hk_3.}
\tag{P4.58}
\end{align}

Physical stage states are
\begin{align}
\boxed{
x_n^{(q)}
=
\mu_x+S_xz_n^{(q)}.
}
\tag{P4.59}
\end{align}

The physical derivative at each stage is
\begin{align}
\boxed{
\widetilde f_\omega^{(q)}
=
\frac{S_x}{\Delta t}
g_\omega(z_n^{(q)},\widetilde v_k).
}
\tag{P4.60}
\end{align}

\section*{14. Simple normalization picture}

There are three computational spaces:
\begin{align}
\boxed{
\text{physical data}
\rightarrow
\text{normalized neural coordinates}
\rightarrow
\text{physical outputs and physics quantities}.
}
\tag{P4.61}
\end{align}

For example, if
\begin{align}
\mu_T=22^\circ{\rm C},
\qquad
s_T=2^\circ{\rm C},
\end{align}
then
\begin{align}
\boxed{
T=24^\circ{\rm C}
\Longrightarrow
z_T=\frac{24-22}{2}=1,
}
\tag{P4.62}
\end{align}
while
\begin{align}
\boxed{
T=20^\circ{\rm C}
\Longrightarrow
z_T=-1.
}
\tag{P4.63}
\end{align}

\section*{15. Where normalization is used}

Before the neural network,
\begin{align}
\boxed{z=S_x^{-1}(x-\mu_x),}
\tag{P4.64}
\end{align}
and
\begin{align}
\boxed{\widetilde v=S_v^{-1}(v-\mu_v).}
\tag{P4.65}
\end{align}

The neural network computes
\begin{align}
\boxed{\frac{dz}{d\tau}=g_\omega(z,\widetilde v).}
\tag{P4.66}
\end{align}

The corresponding physical derivative is
\begin{align}
\boxed{
\widetilde f_\omega
=
\frac{S_x}{\Delta t}g_\omega(z,\widetilde v).
}
\tag{P4.67}
\end{align}

The neural network and RK4 solver therefore operate in normalized coordinates, while the physical RC residual is evaluated in physical units.

\section*{16. Construction and use of \(x\), \(z\), and \(y\)}

For one-zone 1C,
\begin{align}
\boxed{x=T_z.}
\tag{P4.68}
\end{align}

For coupled identity 1C,
\begin{align}
\boxed{
x=
\begin{bmatrix}
T_D\\T_K
\end{bmatrix}.
}
\tag{P4.69}
\end{align}

For one-zone 2C,
\begin{align}
\boxed{
x=
\begin{bmatrix}
T_a\\T_m
\end{bmatrix}.
}
\tag{P4.70}
\end{align}

For coupled 2C,
\begin{align}
\boxed{
x=
\begin{bmatrix}
T_{a,D}\\T_{m,D}\\T_{a,K}\\T_{m,K}
\end{bmatrix}.
}
\tag{P4.71}
\end{align}

The normalized state
\begin{align}
\boxed{z=S_x^{-1}(x-\mu_x)}
\tag{P4.72}
\end{align}
is the internal neural/ODE-solver state.

Whenever physical quantities are required,
\begin{align}
\boxed{x=\mu_x+S_xz.}
\tag{P4.73}
\end{align}

The output is
\begin{align}
\boxed{y=Hx.}
\tag{P4.74}
\end{align}

For one-zone 2C,
\begin{align}
\boxed{
H=
\begin{bmatrix}
1&0
\end{bmatrix},
\qquad
y=T_a.
}
\tag{P4.75}
\end{align}

For coupled 2C,
\begin{align}
\boxed{
H=
\begin{bmatrix}
1&0&0&0\\
0&0&1&0
\end{bmatrix},
}
\tag{P4.76}
\end{align}
so
\begin{align}
\boxed{
y=
\begin{bmatrix}
T_{a,D}\\T_{a,K}
\end{bmatrix}.
}
\tag{P4.77}
\end{align}

Thus
\begin{align}
\boxed{x=\text{full physical state},}
\tag{P4.78}
\end{align}
\begin{align}
\boxed{z=\text{normalized internal neural/solver state},}
\tag{P4.79}
\end{align}
and
\begin{align}
\boxed{y=\text{measured/output portion of }x.}
\tag{P4.80}
\end{align}

\section*{17. Where \(S_y\) is used}

The physical output error is
\begin{align}
e_y=\widehat y-y.
\end{align}

The loss uses
\begin{align}
\boxed{
\widetilde e_y
=
S_y^{-1}(\widehat y-y).
}
\tag{P4.81}
\end{align}

Thus data-loss terms use
\begin{align}
\boxed{
\left\|
S_y^{-1}(\widehat y-y)
\right\|_2^2.
}
\tag{P4.82}
\end{align}

Importantly,
\begin{align}
\boxed{
S_y\text{ does not alter saved physical outputs.}
}
\tag{P4.83}
\end{align}

Predicted temperatures remain in physical units; \(S_y\) only balances their contributions to optimization.

\section*{18. Complete normalization flow}

The forward path is
\begin{align}
\boxed{
x_{\rm physical}
\rightarrow
z
\rightarrow
g_\omega
\rightarrow
z_{\rm next}
\rightarrow
x_{\rm next}
\rightarrow
y_{\rm next}.
}
\tag{P4.84}
\end{align}

Explicitly,
\begin{align}
z_k&=S_x^{-1}(x_k-\mu_x),\\
\widetilde v_k&=S_v^{-1}(v_k-\mu_v),\\
\dot z&=g_\omega(z,\widetilde v),\\
x_{k+1}&=\mu_x+S_xz_{k+1},\\
\widehat y_{k+1}&=Hx_{k+1}.
\end{align}

The data-loss error is
\begin{align}
\boxed{
S_y^{-1}(\widehat y_{k+1}-y_{k+1}).
}
\tag{P4.85}
\end{align}

For physics,
\begin{align}
z&\rightarrow x=\mu_x+S_xz,\\
g_\omega&\rightarrow
\widetilde f_\omega=\frac{S_x}{\Delta t}g_\omega.
\end{align}

Then
\begin{align}
\boxed{
r_{\rm RC}
=
M_\theta\widetilde f_\omega-b_\theta(x,v).
}
\tag{P4.86}
\end{align}

\section*{19. Why RC physics is evaluated in physical units}

The RC residual has consistent power units.

For example,
\begin{align}
\boxed{
C\dot T
:
\frac{J}{K}\frac{K}{s}=W,
}
\tag{P4.87}
\end{align}
and
\begin{align}
\boxed{
\frac{T_o-T}{R}
:
\frac{K}{K/W}=W.
}
\tag{P4.88}
\end{align}

Heat inputs \(Q\) are also in watts. Therefore physics is evaluated after converting the neural derivative and state back to physical units.

\section*{20. Why 2C requires a causal encoder}

For 1C the full state is measured:
\begin{align}
\boxed{x_k=T_{z,k}=y_k.}
\tag{P4.89}
\end{align}

For 2C,
\begin{align}
\boxed{
x_k=
\begin{bmatrix}
T_{a,k}\\T_{m,k}
\end{bmatrix},
}
\tag{P4.90}
\end{align}
but only
\begin{align}
\boxed{y_k=T_{a,k}}
\tag{P4.91}
\end{align}
is measured.

Hence \(T_{m,k}\) must be causally inferred:
\begin{align}
\boxed{
T_{m,k}^{(0)}
=
T_{a,k}
+
\Delta T_m^{\max}
\tanh[e_\psi(\mathcal C_k)].
}
\tag{P4.92}
\end{align}

The context is
\begin{align}
\boxed{
\mathcal C_k
=
\{(y_j,v_j):j=k-L_e+1,\ldots,k\}.
}
\tag{P4.93}
\end{align}

No future observations are allowed.

\section*{21. Simple causal-encoder example}

Suppose
\begin{align}
T_a(09{:}00)=21.6^\circ{\rm C}.
\end{align}

A warming history might be
\begin{align}
21.0,\quad21.2,\quad21.4,\quad21.6^\circ{\rm C},
\end{align}
for which a plausible lagging mass state might be
\begin{align}
T_m(09{:}00)\approx21.1^\circ{\rm C}.
\end{align}

A cooling history could have the same current air temperature while the mass remains warmer, for example
\begin{align}
T_m(09{:}00)\approx22.2^\circ{\rm C}.
\end{align}

Thus the same instantaneous air temperature does not uniquely determine the hidden thermal mass state.

The encoder therefore uses
\begin{align}
\boxed{
T_m(t_k)
=
E_\psi
\left(
T_a(t_{k-L_e+1:k}),
v(t_{k-L_e+1:k})
\right).
}
\tag{P4.94}
\end{align}

\section*{22. Where the \(T_m\) used in mass residuals comes from}

At rollout initialization,
\begin{align}
\boxed{T_{m,k}=E_\psi(\mathcal C_k).}
\tag{P4.95}
\end{align}

After initialization, the encoder is not repeatedly called within the same rollout.

The PINODE predicts \(\widetilde{\dot T}_m\), and RK4 propagates the hidden state:
\begin{align}
\boxed{
\text{causal encoder}
\rightarrow
T_{m,k}
\rightarrow
\text{PINODE integration}
\rightarrow
T_{m,k+\delta}
\rightarrow
T_{m,k+1}
\rightarrow\cdots.
}
\tag{P4.96}
\end{align}

Therefore the \(T_m\) in each mass residual is the current dynamically propagated hidden model state.

\section*{23. Physics loss also trains the encoder}

For 2C,
\begin{align}
\boxed{
\mathcal J_{\rm PINODE}
=
\mathcal J_{\rm PINODE}(\omega,\psi,\theta_{RC}).
}
\tag{P4.97}
\end{align}

The encoder gradient is
\begin{align}
\boxed{
\frac{\partial\mathcal J}{\partial\psi}
=
\lambda_y
\frac{\partial\mathcal L_{\rm data}}{\partial\psi}
+
\lambda_f
\frac{\partial\mathcal L_{\rm phys}}{\partial\psi}
+
\lambda_{\rm wd}
\frac{\partial\mathcal L_{\rm reg}}{\partial\psi}.
}
\tag{P4.98}
\end{align}

\section*{24. All-to-One 1C PINODE}

Define
\begin{align}
\boxed{
Q_{c,A}
=
Q_{ZIC,A}+Q_{Sol1,A}+Q_{AC,A},
}
\tag{P4.99}
\end{align}
and
\begin{align}
\boxed{
Q_{r,A}
=
Q_{ZIR,A}+Q_{Sol2,A}.
}
\tag{P4.100}
\end{align}

The raw derivative is
\begin{align}
\boxed{
\widetilde{\dot T}_A
=
\widetilde f_{\omega,A}^{1C}(T_A,v_A).
}
\tag{P4.101}
\end{align}

The residual is
\begin{align}
\boxed{
\begin{aligned}
r_A^{1C}
={}&
C_A\widetilde{\dot T}_A
-
\frac{T_o-T_A}{R_{Ao}}
-
Q_{c,A}
-
\eta_{r,A}Q_{r,A}.
\end{aligned}
}
\tag{P4.102}
\end{align}

The normalized contribution is
\begin{align}
\boxed{
\ell_{\rm phys,A}^{1C}
=
\left(\frac{r_A^{1C}}{q_A^\star}\right)^2.
}
\tag{P4.103}
\end{align}

\section*{25. All-to-One 2C PINODE}

State and raw derivative are
\begin{align}
\boxed{
x_A=
\begin{bmatrix}
T_{a,A}\\T_{m,A}
\end{bmatrix},
\qquad
\widetilde f_{\omega,A}^{2C}
=
\begin{bmatrix}
\widetilde{\dot T}_{a,A}\\
\widetilde{\dot T}_{m,A}
\end{bmatrix}.
}
\tag{P4.104}
\end{align}

Air residual:
\begin{align}
\boxed{
\begin{aligned}
r_{a,A}
={}&
C_{a,A}\widetilde{\dot T}_{a,A}
-
\frac{T_o-T_{a,A}}{R_{Ao}}
-
\frac{T_{m,A}-T_{a,A}}{R_{Am}}
\\
&-
Q_{c,A}
-
(1-\eta_{r,A})Q_{r,A}.
\end{aligned}
}
\tag{P4.105}
\end{align}

Mass residual:
\begin{align}
\boxed{
\begin{aligned}
r_{m,A}
={}&
C_{m,A}\widetilde{\dot T}_{m,A}
-
\frac{T_{a,A}-T_{m,A}}{R_{Am}}
-
\eta_{r,A}Q_{r,A}.
\end{aligned}
}
\tag{P4.106}
\end{align}

Residual vector:
\begin{align}
\boxed{
r_A^{2C}
=
\begin{bmatrix}
r_{a,A}\\r_{m,A}
\end{bmatrix}.
}
\tag{P4.107}
\end{align}

The total-energy diagnostic is
\begin{align}
\boxed{
\begin{aligned}
r_{\Sigma,A}
=
{}&
C_{a,A}\widetilde{\dot T}_{a,A}
+
C_{m,A}\widetilde{\dot T}_{m,A}
-
\frac{T_o-T_{a,A}}{R_{Ao}}
-
Q_{c,A}
-
Q_{r,A}.
\end{aligned}
}
\tag{P4.108}
\end{align}

\section*{26. Identity / IND 1C PINODE}

Dining:
\begin{align}
\boxed{
r_D
=
C_D\widetilde{\dot T}_D
-
\frac{T_o-T_D}{R_{Do}}
-
Q_{c,D}
-
\eta_{r,D}Q_{r,D}.
}
\tag{P4.109}
\end{align}

Kitchen:
\begin{align}
\boxed{
r_K
=
C_K\widetilde{\dot T}_K
-
\frac{T_o-T_K}{R_{Ko}}
-
Q_{c,K}
-
\eta_{r,K}Q_{r,K}.
}
\tag{P4.110}
\end{align}

For Kitchen,
\begin{align}
\boxed{
Q_{c,K}=Q_{ZIC,K}+Q_{AC,K},
\qquad
Q_{r,K}=Q_{ZIR,K}.
}
\tag{P4.111}
\end{align}

No \(R_{DK}\) exists in IND.

\section*{27. Identity / IND 2C PINODE}

Dining:
\begin{align}
\boxed{
\begin{aligned}
r_{a,D}
={}&
C_{a,D}\widetilde{\dot T}_{a,D}
-
\frac{T_o-T_{a,D}}{R_{Do}}
-
\frac{T_{m,D}-T_{a,D}}{R_{Dm}}
-
Q_{c,D}
-
(1-\eta_{r,D})Q_{r,D},
\end{aligned}
}
\tag{P4.112}
\end{align}
\begin{align}
\boxed{
r_{m,D}
=
C_{m,D}\widetilde{\dot T}_{m,D}
-
\frac{T_{a,D}-T_{m,D}}{R_{Dm}}
-
\eta_{r,D}Q_{r,D}.
}
\tag{P4.113}
\end{align}

Kitchen:
\begin{align}
\boxed{
\begin{aligned}
r_{a,K}
={}&
C_{a,K}\widetilde{\dot T}_{a,K}
-
\frac{T_o-T_{a,K}}{R_{Ko}}
-
\frac{T_{m,K}-T_{a,K}}{R_{Km}}
-
Q_{c,K}
-
(1-\eta_{r,K})Q_{r,K},
\end{aligned}
}
\tag{P4.114}
\end{align}
\begin{align}
\boxed{
r_{m,K}
=
C_{m,K}\widetilde{\dot T}_{m,K}
-
\frac{T_{a,K}-T_{m,K}}{R_{Km}}
-
\eta_{r,K}Q_{r,K}.
}
\tag{P4.115}
\end{align}

\section*{28. Identity / DEP1 1C PINODE}

State and derivative:
\begin{align}
\boxed{
x=
\begin{bmatrix}
T_D\\T_K
\end{bmatrix},
\qquad
\widetilde f_{\omega,\rm DEP1}^{1C}
=
\begin{bmatrix}
\widetilde{\dot T}_D\\
\widetilde{\dot T}_K
\end{bmatrix}.
}
\tag{P4.116}
\end{align}

Dining:
\begin{align}
\boxed{
r_D
=
C_D\widetilde{\dot T}_D
-
\frac{T_o-T_D}{R_{Do}}
-
\frac{T_K-T_D}{R_{DK}}
-
Q_{c,D}
-
\eta_{r,D}Q_{r,D}.
}
\tag{P4.117}
\end{align}

Kitchen:
\begin{align}
\boxed{
r_K
=
C_K\widetilde{\dot T}_K
-
\frac{T_o-T_K}{R_{Ko}}
-
\frac{T_D-T_K}{R_{DK}}
-
Q_{c,K}
-
\eta_{r,K}Q_{r,K}.
}
\tag{P4.118}
\end{align}

Residual vector:
\begin{align}
\boxed{
r_{\rm DEP1}^{1C}
=
\begin{bmatrix}
r_D\\r_K
\end{bmatrix}.
}
\tag{P4.119}
\end{align}

The coupling is reciprocal:
\begin{align}
\boxed{R_{DK}=R_{KD}.}
\tag{P4.120}
\end{align}

\section*{29. Identity / DEP1 2C PINODE}

State:
\begin{align}
\boxed{
x=
\begin{bmatrix}
T_{a,D}\\T_{m,D}\\T_{a,K}\\T_{m,K}
\end{bmatrix}.
}
\tag{P4.121}
\end{align}

Dining air:
\begin{align}
\boxed{
\begin{aligned}
r_{a,D}
={}&
C_{a,D}\widetilde{\dot T}_{a,D}
-
\frac{T_o-T_{a,D}}{R_{Do}}
-
\frac{T_{m,D}-T_{a,D}}{R_{Dm}}
-
\frac{T_{a,K}-T_{a,D}}{R_{DK}}
\\
&-
Q_{c,D}
-
(1-\eta_{r,D})Q_{r,D}.
\end{aligned}
}
\tag{P4.122}
\end{align}

Dining mass:
\begin{align}
\boxed{
r_{m,D}
=
C_{m,D}\widetilde{\dot T}_{m,D}
-
\frac{T_{a,D}-T_{m,D}}{R_{Dm}}
-
\eta_{r,D}Q_{r,D}.
}
\tag{P4.123}
\end{align}

Kitchen air:
\begin{align}
\boxed{
\begin{aligned}
r_{a,K}
={}&
C_{a,K}\widetilde{\dot T}_{a,K}
-
\frac{T_o-T_{a,K}}{R_{Ko}}
-
\frac{T_{m,K}-T_{a,K}}{R_{Km}}
-
\frac{T_{a,D}-T_{a,K}}{R_{DK}}
\\
&-
Q_{c,K}
-
(1-\eta_{r,K})Q_{r,K}.
\end{aligned}
}
\tag{P4.124}
\end{align}

Kitchen mass:
\begin{align}
\boxed{
r_{m,K}
=
C_{m,K}\widetilde{\dot T}_{m,K}
-
\frac{T_{a,K}-T_{m,K}}{R_{Km}}
-
\eta_{r,K}Q_{r,K}.
}
\tag{P4.125}
\end{align}

Residual vector:
\begin{align}
\boxed{
r_{\rm DEP1}^{2C}
=
\begin{bmatrix}
r_{a,D}\\r_{m,D}\\r_{a,K}\\r_{m,K}
\end{bmatrix}.
}
\tag{P4.126}
\end{align}

Zone-total residuals are
\begin{align}
\boxed{
\begin{aligned}
r_{\Sigma,D}
={}&
C_{a,D}\widetilde{\dot T}_{a,D}
+
C_{m,D}\widetilde{\dot T}_{m,D}
-
\frac{T_o-T_{a,D}}{R_{Do}}
-
\frac{T_{a,K}-T_{a,D}}{R_{DK}}
-
Q_{c,D}
-
Q_{r,D},
\end{aligned}
}
\tag{P4.127}
\end{align}
and
\begin{align}
\boxed{
\begin{aligned}
r_{\Sigma,K}
={}&
C_{a,K}\widetilde{\dot T}_{a,K}
+
C_{m,K}\widetilde{\dot T}_{m,K}
-
\frac{T_o-T_{a,K}}{R_{Ko}}
-
\frac{T_{a,D}-T_{a,K}}{R_{DK}}
-
Q_{c,K}
-
Q_{r,K}.
\end{aligned}
}
\tag{P4.128}
\end{align}

\section*{30. Identity / DEP2 1C PINODE}

Define
\begin{align}
\boxed{
\bar Q_c^{nh}
=
Q_{ZIC,A}+Q_{Sol1,A},
}
\tag{P4.129}
\end{align}
and
\begin{align}
\boxed{
\bar Q_r
=
Q_{ZIR,A}+Q_{Sol2,A}.
}
\tag{P4.130}
\end{align}

Dining:
\begin{align}
\boxed{
\begin{aligned}
r_D
={}&
C_D\widetilde{\dot T}_D
-
\frac{T_o-T_D}{R_{Do}}
-
\frac{T_K-T_D}{R_{DK}}
-
Q_{AC,D}
-
\lambda_{c,D}\bar Q_c^{nh}
-
\eta_{r,D}\lambda_{r,D}\bar Q_r.
\end{aligned}
}
\tag{P4.131}
\end{align}

Kitchen:
\begin{align}
\boxed{
\begin{aligned}
r_K
={}&
C_K\widetilde{\dot T}_K
-
\frac{T_o-T_K}{R_{Ko}}
-
\frac{T_D-T_K}{R_{DK}}
-
Q_{AC,K}
-
\lambda_{c,K}\bar Q_c^{nh}
-
\eta_{r,K}\lambda_{r,K}\bar Q_r.
\end{aligned}
}
\tag{P4.132}
\end{align}

For the primary DEP2 1C model,
\begin{align}
\boxed{\eta_{r,D}=\eta_{r,K}=1.}
\tag{P4.133}
\end{align}

\section*{31. Identity / DEP2 2C PINODE}

Dining air:
\begin{align}
\boxed{
\begin{aligned}
r_{a,D}
={}&
C_{a,D}\widetilde{\dot T}_{a,D}
-
\frac{T_o-T_{a,D}}{R_{Do}}
-
\frac{T_{m,D}-T_{a,D}}{R_{Dm}}
-
\frac{T_{a,K}-T_{a,D}}{R_{DK}}
-
Q_{AC,D}
-
\lambda_{c,D}\bar Q_c^{nh}
-
(1-\eta_{r,D})\lambda_{r,D}\bar Q_r.
\end{aligned}
}
\tag{P4.134}
\end{align}

Dining mass:
\begin{align}
\boxed{
r_{m,D}
=
C_{m,D}\widetilde{\dot T}_{m,D}
-
\frac{T_{a,D}-T_{m,D}}{R_{Dm}}
-
\eta_{r,D}\lambda_{r,D}\bar Q_r.
}
\tag{P4.135}
\end{align}

Kitchen air:
\begin{align}
\boxed{
\begin{aligned}
r_{a,K}
={}&
C_{a,K}\widetilde{\dot T}_{a,K}
-
\frac{T_o-T_{a,K}}{R_{Ko}}
-
\frac{T_{m,K}-T_{a,K}}{R_{Km}}
-
\frac{T_{a,D}-T_{a,K}}{R_{DK}}
-
Q_{AC,K}
-
\lambda_{c,K}\bar Q_c^{nh}
-
(1-\eta_{r,K})\lambda_{r,K}\bar Q_r.
\end{aligned}
}
\tag{P4.136}
\end{align}

Kitchen mass:
\begin{align}
\boxed{
r_{m,K}
=
C_{m,K}\widetilde{\dot T}_{m,K}
-
\frac{T_{a,K}-T_{m,K}}{R_{Km}}
-
\eta_{r,K}\lambda_{r,K}\bar Q_r.
}
\tag{P4.137}
\end{align}

Residual vector:
\begin{align}
\boxed{
r_{\rm DEP2}^{2C}
=
\begin{bmatrix}
r_{a,D}\\r_{m,D}\\r_{a,K}\\r_{m,K}
\end{bmatrix}.
}
\tag{P4.138}
\end{align}

The zone-total residuals are
\begin{align}
\boxed{
\begin{aligned}
r_{\Sigma,D}
={}&
C_{a,D}\widetilde{\dot T}_{a,D}
+
C_{m,D}\widetilde{\dot T}_{m,D}
-
\frac{T_o-T_{a,D}}{R_{Do}}
-
\frac{T_{a,K}-T_{a,D}}{R_{DK}}
-
Q_{AC,D}
-
\lambda_{c,D}\bar Q_c^{nh}
-
\lambda_{r,D}\bar Q_r,
\end{aligned}
}
\tag{P4.139}
\end{align}
and
\begin{align}
\boxed{
\begin{aligned}
r_{\Sigma,K}
={}&
C_{a,K}\widetilde{\dot T}_{a,K}
+
C_{m,K}\widetilde{\dot T}_{m,K}
-
\frac{T_o-T_{a,K}}{R_{Ko}}
-
\frac{T_{a,D}-T_{a,K}}{R_{DK}}
-
Q_{AC,K}
-
\lambda_{c,K}\bar Q_c^{nh}
-
\lambda_{r,K}\bar Q_r.
\end{aligned}
}
\tag{P4.140}
\end{align}

\section*{32. Physics-residual normalization}

For zone \(i\), define a training-only heat scale
\begin{align}
\boxed{
q_i^\star
=
\sqrt{
\frac{1}{N_{\rm tr}}
\sum_{k\in\mathcal I_{\rm tr}}
\left(
Q_{c,i,k}^2+Q_{r,i,k}^2
\right)
}
+
\epsilon_q.
}
\tag{P4.141}
\end{align}

Then
\begin{align}
\boxed{
\widetilde r_i=\frac{r_i}{q_i^\star}.
}
\tag{P4.142}
\end{align}

For 2C,
\begin{align}
\boxed{
\ell_{\rm phys,i}^{2C}
=
\frac{r_{a,i}^2+r_{m,i}^2}{(q_i^\star)^2}.
}
\tag{P4.143}
\end{align}

The same zone heat scale is used for both air and mass equations.

\section*{33. Data rollout loss}

For rollout start \(k\),
\begin{align}
\boxed{
\mathcal L_{\rm data}^{(k)}
=
\frac{1}{N_r}
\sum_{\ell=1}^{N_r}
\left\|
S_y^{-1}
\left(
\widehat y_{k+\ell|k}-y_{k+\ell}
\right)
\right\|_2^2.
}
\tag{P4.144}
\end{align}

Across rollout windows,
\begin{align}
\boxed{
\mathcal L_{\rm data}
=
\frac{1}{|\mathcal W_{\rm tr}|}
\sum_{k\in\mathcal W_{\rm tr}}
\mathcal L_{\rm data}^{(k)}.
}
\tag{P4.145}
\end{align}

Predicted states are recursively fed back within each rollout.

\section*{34. Complete base-PINODE objective}

The core objective is
\begin{align}
\boxed{
\mathcal J_{\rm PINODE}
=
\lambda_y\mathcal L_{\rm data}
+
\lambda_f\mathcal L_{\rm phys}.
}
\tag{P4.146}
\end{align}

With weight decay,
\begin{align}
\boxed{
\mathcal J_{\rm PINODE}
=
\lambda_y\mathcal L_{\rm data}
+
\lambda_f\mathcal L_{\rm phys}
+
\lambda_{\rm wd}\mathcal L_{\rm reg}.
}
\tag{P4.147}
\end{align}

For 2C,
\begin{align}
\boxed{
\mathcal L_{\rm reg}
=
\|\omega\|_2^2+\|\psi\|_2^2.
}
\tag{P4.148}
\end{align}

For 1C the encoder term is absent.

If
\begin{align}
\lambda_f=0,
\end{align}
then
\begin{align}
\boxed{\text{base PINODE reduces to NODE.}}
\tag{P4.149}
\end{align}

For finite \(\lambda_f\),
\begin{align}
\boxed{r_{\rm RC}=0\text{ is not guaranteed.}}
\tag{P4.150}
\end{align}

\section*{35. Training stages}

Stage A is a NODE warm start using
\begin{align}
\boxed{\mathcal L_{\rm data}.}
\tag{P4.151}
\end{align}

Stage B jointly optimizes neural and physical parameters:
\begin{align}
\boxed{
\mathcal J
=
\lambda_y\mathcal L_{\rm data}
+
\lambda_f\mathcal L_{\rm phys}.
}
\tag{P4.152}
\end{align}

Stage C performs full joint refinement until validation rollout performance stops improving.

\section*{36. Validation criterion}

For validation segment \(s\),
\begin{align}
\boxed{
\mathcal E_s^{val}
=
\frac{1}{N_{\rm val}}
\sum_{\ell=1}^{N_{\rm val}}
\left\|
S_y^{-1}
(\widehat y_{s,\ell}-y_{s,\ell})
\right\|_2^2.
}
\tag{P4.153}
\end{align}

Across validation segments,
\begin{align}
\boxed{
\mathcal E_{\rm val}
=
\frac{1}{N_{\rm seg,val}}
\sum_s\mathcal E_s^{val}.
}
\tag{P4.154}
\end{align}

Checkpoint selection is
\begin{align}
\boxed{
(\widehat\omega,\widehat\psi,\widehat\theta)
=
\arg\min_{\rm checkpoint}
\mathcal E_{\rm val}.
}
\tag{P4.155}
\end{align}

Physics residual is logged but is not the sole checkpoint selector.

\section*{37. Locked hyperparameter-selection protocol}

Hyperparameter tuning is first performed on a smaller representative subset
\begin{align}
\boxed{
\mathcal D_{\rm tune}\subset\mathcal D_{\rm train}.
}
\tag{P4.156}
\end{align}

The tuning subset is constructed only from training data and should preserve variation in heating/cooling conditions, occupied/unoccupied operation, QAC, solar gains, and outdoor temperature.

For candidate hyperparameters \(h\),
\begin{align}
\boxed{
h^\star
=
\arg\min_h
\mathcal E_{\rm tune,val}(h).
}
\tag{P4.157}
\end{align}

Then
\begin{align}
\boxed{h=h^\star}
\tag{P4.158}
\end{align}
is frozen.

The final model is then trained on the complete training data:
\begin{align}
\boxed{
\widehat\Theta
=
\arg\min_\Theta
\mathcal J
(
\Theta;
\mathcal D_{\rm train},
h^\star
).
}
\tag{P4.159}
\end{align}

The test set is untouched until final evaluation.

This reduced-data tuning followed by frozen-hyperparameter full-data training is used for all four thermal-model methods.

\section*{38. Hyperparameters included in tuning}

The PINODE tuning set includes
\begin{align}
\boxed{
\mathcal H_{\rm PINODE}
=
\{
L_{\rm hid},
\bm n_{\rm hid},
\varphi,
N_r,
L_e,
N_s,
\lambda_f,
\lambda_{\rm wd},
\text{learning rate},
\text{batch size},
\text{optimizer},
\Delta T_m^{\max}
\}.
}
\tag{P4.160}
\end{align}

In particular,
\begin{align}
\boxed{
L_{\rm hid}
=
\text{number of hidden layers},
}
\tag{P4.161}
\end{align}
and
\begin{align}
\boxed{
\bm n_{\rm hid}
=
\text{number of neurons in the hidden layers}
}
\tag{P4.162}
\end{align}
are explicit hyperparameters.

For 1C, \(L_e\) and \(\Delta T_m^{\max}\) do not apply.

Once selected, all hyperparameters are frozen for final full-data training.

\section*{39. Sim1}

For 1C,
\begin{align}
\boxed{x_k=y_k.}
\tag{P4.163}
\end{align}

Then
\begin{align}
\boxed{
\widehat x_{k+1}
=
\Phi_{\widehat\omega}
(x_k,v_k^{test};N_s).
}
\tag{P4.164}
\end{align}

For 2C,
\begin{align}
\boxed{
x_k=[y_k,E_{\widehat\psi}(\mathcal C_k)].
}
\tag{P4.165}
\end{align}

The measured air state is reset each step and the hidden state is reconstructed causally.

\section*{40. Sim2}

Initialize once:
\begin{align}
\boxed{\widehat x_0=y_0}
\tag{P4.166}
\end{align}
for 1C, or
\begin{align}
\boxed{
\widehat x_0=[y_0,E_{\widehat\psi}(\mathcal C_0)]
}
\tag{P4.167}
\end{align}
for 2C.

Then recursively
\begin{align}
\boxed{
\widehat x_{k+1}
=
\Phi_{\widehat\omega}
(\widehat x_k,v_k^{test};N_s).
}
\tag{P4.168}
\end{align}

Recorded test QAC and disturbances remain external forcing.

\section*{41. Sim3}

The local thermostat produces mass flow:
\begin{align}
\boxed{
\dot m_k
=
\pi_{\rm thermostat}
(\widehat T_{z,k},T_{\rm set,k}).
}
\tag{P4.169}
\end{align}

The Phase C QAC model produces
\begin{align}
\boxed{
Q_{AC,k}
=
G_{QAC}
(\dot m_k,\widehat T_{z,k},T_{sa,k},\ldots).
}
\tag{P4.170}
\end{align}

The PINODE integrates
\begin{align}
\boxed{
\widehat x_{k+1}
=
\Phi_{\widehat\omega}
(\widehat x_k,Q_{AC,k},d_k^{test};N_s).
}
\tag{P4.171}
\end{align}

Only disturbances are replayed from test data.

\section*{42. Locked PHVAC contract for Sim1/Sim2/Sim3}

For every thermal-model method and every Sim1/Sim2/Sim3 evaluation, PHVAC is computed and saved using the common Phase C PHVAC model.

The common mapping is
\begin{align}
\boxed{
\widehat P_{\rm HVAC,k}
=
G_{\rm PHVAC}(Q_{AC,k},\ldots).
}
\tag{P4.172}
\end{align}

For Sim1,
\begin{align}
\boxed{Q_{AC,k}=Q_{AC,k}^{test}.}
\tag{P4.173}
\end{align}

For Sim2,
\begin{align}
\boxed{Q_{AC,k}=Q_{AC,k}^{test}.}
\tag{P4.174}
\end{align}

For Sim3,
\begin{align}
\boxed{
Q_{AC,k}
=
G_{QAC}
(\dot m_k,\widehat T_{z,k},T_{sa,k},\ldots).
}
\tag{P4.175}
\end{align}

Then in every simulation,
\begin{align}
\boxed{
\widehat P_{\rm HVAC,k}
=
G_{\rm PHVAC}(Q_{AC,k},\ldots).
}
\tag{P4.176}
\end{align}

For Sim1 and Sim2, PHVAC may be identical across thermal models when the same recorded QAC trajectory is replayed. It is still saved to preserve a common output contract.

Required outputs for Sim1 and Sim2 include
\begin{align}
\boxed{
\widehat T_z(t),
\quad
Q_{AC}^{test}(t),
\quad
\widehat P_{\rm HVAC}(t).
}
\tag{P4.177}
\end{align}

Required outputs for Sim3 include
\begin{align}
\boxed{
\widehat T_z(t),
\quad
\dot m(t),
\quad
Q_{AC}(t),
\quad
\widehat P_{\rm HVAC}(t).
}
\tag{P4.178}
\end{align}

\section*{43. Residual diagnostics}

At minimum record
\begin{align}
\boxed{\mathcal L_{\rm data},}
\end{align}
\begin{align}
\boxed{\mathcal L_{\rm phys},}
\end{align}
\begin{align}
\boxed{\mathcal J_{\rm PINODE},}
\end{align}
and validation rollout error.

For coupled 2C also record
\begin{align}
\boxed{
\operatorname{RMSE}(r_{a,D}),
\quad
\operatorname{RMSE}(r_{m,D}),
}
\tag{P4.179}
\end{align}
\begin{align}
\boxed{
\operatorname{RMSE}(r_{a,K}),
\quad
\operatorname{RMSE}(r_{m,K}),
}
\tag{P4.180}
\end{align}
and total-energy residuals
\begin{align}
\boxed{r_{\Sigma,D},\quad r_{\Sigma,K}.}
\tag{P4.181}
\end{align}

Define the raw neural physics violation
\begin{align}
\boxed{
\rho
=
r_{\rm RC}
=
M_\theta\widetilde f_\omega-b_\theta.
}
\tag{P4.182}
\end{align}

\section*{44. Base PINODE versus other methods}

Inverse PINN deploys
\begin{align}
\boxed{
f_{RC}(x,v;\widehat\theta).
}
\tag{P4.183}
\end{align}

Base PINODE deploys
\begin{align}
\boxed{
\widetilde f_{\widehat\omega}(x,v).
}
\tag{P4.184}
\end{align}

NODE uses
\begin{align}
\boxed{
\mathcal J_{\rm NODE}=\mathcal L_{\rm data}.
}
\tag{P4.185}
\end{align}

Base PINODE uses
\begin{align}
\boxed{
\mathcal J_{\rm PINODE}
=
\mathcal L_{\rm data}
+
\lambda_f\mathcal L_{\rm phys}
}
\tag{P4.186}
\end{align}
up to common weights and regularization.

Base PINODE satisfies RC equations only softly:
\begin{align}
\boxed{r_{\rm RC}\approx0.}
\tag{P4.187}
\end{align}

EBP will instead construct a projected derivative satisfying selected energy constraints
\begin{align}
\boxed{Af_P=b}
\tag{P4.188}
\end{align}
by construction.

Thus
\begin{align}
\boxed{\text{soft physics}\neq\text{hard physics}.}
\tag{P4.189}
\end{align}

\section*{45. Complete base-PINODE algorithm}

Choose
\begin{align}
\boxed{
\mathcal A\in\{A,\mathrm{IND},\mathrm{DEP1},\mathrm{DEP2}\},
}
\end{align}
and choose
\begin{align}
\boxed{1C\quad\text{or}\quad2C.}
\end{align}

Load the corresponding Phase-D \texttt{l1\_h1} data, build contiguous segments, compute training-only normalization and \(q_i^\star\), construct rollout windows of length \(N_r\), and for 2C construct causal contexts of length \(L_e\).

Evaluate
\begin{align}
\boxed{
\widetilde f_\omega
=
\frac{S_x}{\Delta t}g_\omega(z,\widetilde v).
}
\tag{P4.190}
\end{align}

Integrate using differentiable RK4 with \(N_s\) substeps per sampled interval.

At every RK4 stage compute
\begin{align}
\boxed{
r_{\rm RC}
=
M_\theta\widetilde f_\omega-b_\theta.
}
\tag{P4.191}
\end{align}

Optimize
\begin{align}
\boxed{
\min_{\omega,\psi,\theta_{RC}}
\left[
\lambda_y\mathcal L_{\rm data}
+
\lambda_f\mathcal L_{\rm phys}
+
\lambda_{\rm wd}\mathcal L_{\rm reg}
\right].
}
\tag{P4.192}
\end{align}

Hyperparameters are selected first on the smaller representative tuning dataset, including network hidden-layer count and neurons per hidden layer. Hyperparameters are then frozen. Final full-data training is performed on the complete training set, with validation used for checkpoint selection. The test set remains untouched until final Sim1/Sim2/Sim3 evaluation.

The deployed derivative remains
\begin{align}
\boxed{
\dot x
=
\widetilde f_{\widehat\omega}(x,v).
}
\tag{P4.193}
\end{align}

No RC correction or projection is applied during base-PINODE simulation.

\section*{Part 4 conclusion}

The base PINODE is
\begin{align}
\boxed{
\frac{dz}{d\tau}=g_\omega(z,\widetilde v)
}
\end{align}
with physical derivative
\begin{align}
\boxed{
\widetilde f_\omega
=
\frac{S_x}{\Delta t}g_\omega(z,\widetilde v)
}
\end{align}
and full soft RC residual
\begin{align}
\boxed{
r_{\rm RC}
=
M_\theta\widetilde f_\omega-b_\theta.
}
\end{align}

Training minimizes
\begin{align}
\boxed{
\mathcal J_{\rm PINODE}
=
\lambda_y\mathcal L_{\rm data}
+
\lambda_f\mathcal L_{\rm phys}
+
\lambda_{\rm wd}\mathcal L_{\rm reg}.
}
\end{align}

The locked workflow is
\begin{align}
\boxed{\text{representative reduced-data hyperparameter tuning}}
\end{align}
followed by
\begin{align}
\boxed{\text{hyperparameter freezing}}
\end{align}
followed by
\begin{align}
\boxed{\text{full training on complete training data}}
\end{align}
followed by
\begin{align}
\boxed{\text{Sim1/Sim2/Sim3 evaluation on untouched test data}.}
\end{align}

The number of hidden layers and neurons per hidden layer are explicitly tuned in the reduced-data tuning stage and then frozen.

For every method and every Sim1/Sim2/Sim3 run,
\begin{align}
\boxed{
\widehat P_{\rm HVAC}
=
G_{\rm PHVAC}(Q_{AC},\ldots)
}
\end{align}
is computed using the common Phase C PHVAC model and saved.

Thus
\begin{align}
\boxed{\text{NODE}=\text{data-only neural ODE},}
\end{align}
\begin{align}
\boxed{\text{Base PINODE}=\text{neural ODE + soft full-RC residual},}
\end{align}
and Part 5 will define
\begin{align}
\boxed{\text{EBP-PINODE}=\text{neural ODE + exact energy-balance projection}.}
\end{align}

\end{document}
